\documentclass[11pt,a4paper]{article}
\usepackage[margin=25mm]{geometry}
\usepackage{amsmath,amssymb,amsthm,hyperref}
\setlength{\emergencystretch}{2em}
\newtheorem{lemma}{Lemma}
\newtheorem{theorem}[lemma]{Theorem}
\newcommand{\gen}[1]{\langle#1\rangle}
\title{The small anchored imprimitive endpoint}
\author{Human Compression project: bounded L044 proof component}
\date{}
\begin{document}\maketitle
The accepted anchored imprimitive bound covers degree at least twelve.
This note closes its smaller degrees, using the same actual-orbit
forest and the already paid maximal symmetric-string interfaces.
It supplies only the anchored replacement needed in L044's supported
interval induction; it does not restore the false unrestricted
reading of its Proposition~5.13. An endpoint anchor is a point moved
by that endpoint generator and fixed by every other generator.

\begin{lemma}[Small symmetric endpoint test]\label{lem:small}
Let $(h_0,\ldots,h_{b-2})$ be a connected string C-marking of
$S_b$ of length $b-1$, $2\leq b\leq5$. Put
$Q=\gen{h_1,\ldots,h_{b-2}}$ and
$F=\gen{h_0,\ldots,h_{b-3}}$. If $F$ fixes a point $p$, then
\[
                 C_{S_b}(Q)\cap (S_b)_p=1.
\]
Moreover an even-weight binary vector constant on the $Q$-orbits
and zero at $p$ is zero.
\end{lemma}
\begin{proof}
The paid degree-five maximal-string certificate gives the ordinary
adjacent-transposition marking. For $b=2,3$ this is elementary:
distinct involutions of $S_3$ are transpositions. For $b=4$, the
accepted elementary classification gives the ordinary path and,
up to reversal and relabelling,
\[
                  ((12)(34),(23),(34)).
\]
Here is the short classification argument. Both commuting end
involutions cannot be double transpositions: adjoining one
involution to the normal four-group does not generate $S_4$.
For two disjoint single transpositions, the middle must be a
single transposition joining their pairs; a double transposition
preserves that pair partition. For one single and one double
end, they share one factor, and again the middle must join
their pairs by a single transposition. These give exactly the
two forms just displayed.

For $b=2$, a point stabilizer is trivial and an even binary
vector zero at one of the two points is zero. For $b\geq3$ in
an ordinary path labelled $(12),(23),\ldots$, the only
$F$-fixed point is $p=b$. The group $Q=S_{\{2,\ldots,b\}}$
has trivial centralizer fixing $p$: for $b\geq4$ its centralizer
in $S_b$ is trivial; for $b=3$ the possible nonidentity element
swaps $p$. Its large
orbit contains $p$, so a binary vector as stated can be nonzero
only at point~1, which has odd weight.

In the displayed exceptional order, $F$ is transitive $D_8$,
so there is no $p$. In the reversed order, $F$ is a point
stabilizer $S_3$, whereas $Q$ is transitive $D_8$. An element
centralizing a transitive group and fixing a point fixes every
point. A $Q$-constant vector zero at one point is likewise zero.
\end{proof}

\begin{theorem}[Small anchored imprimitive bound]
Let $H\leq A_N$, $4\leq N\leq11$, be transitive and imprimitive,
with a connected string C-marking of rank $d$, all maximal
omissions intransitive, and an endpoint anchor. Then
\[
                         d\leq\left\lfloor\frac{N-2}{2}\right\rfloor.
\]
\end{theorem}
\begin{proof}
Choose blocks of size $a$ and number $b$. Let $M$ be a minimal
distinguished subfamily generating the whole block action, and
let $c$ be its actual number of point orbits. The accepted forest
argument and Whiston bound give
\[
                 d-|M|\leq c-1,\qquad c\leq a,\qquad |M|\leq b-1.
\]
Thus $d\leq a+b-2$. If $a,b\geq3$, the only case here is
$a=b=3$. Equality $d=4$ would force $|M|=2,c=3$ and block
image $S_3$. Its three orbits each meet each block once, so
$\gen M$ is faithfully the diagonal $S_3$ on three copies.
A block transposition lifts to three point transpositions,
contradicting evenness. Hence $d\leq3$ as required.

Suppose $b=2$. Reverse the marking so an anchor $p$ is moved
only by its first generator $\tau$. Every other generator
fixes $p$, hence preserves both blocks; transitivity forces
$\tau$ to interchange them. It is a product of $a$
transpositions, so $a$ is even. Only $a=2,4$ remain.
To exclude $d=a$, normalize the two blocks as row pairs
$1,\ldots,a$ and $1',\ldots,a'$ with $p=1$ and
$\tau=\prod_j(j\ j')$. Put $x=s_1$.

For $a=2$, $x$ fixes~1 and must move~$1'$ to connect its row
pair to the other pair. Thus $x=(1'\ 2')$ is odd.
For $a=4$, the last two generators centralize $\tau$, fix
the anchor row, and are independent. Their diagonal row action
is the full $S_3$, so write them
\[
       h_0=(23)(2'3'),\qquad h_1=(34)(3'4').
\]
The generator $x$ centralizes $h_1$, fixes~1 and preserves
the two blocks. It must move~$1'$; consequently
\[
             x=(1'2')(34)^\epsilon(3'4')^\delta,
             \qquad \epsilon+\delta\equiv1\pmod2.
\]
But $\gen{\tau,x,h_0}$ is transitive: $\tau$ joins the two
points of each row, while the other edges connect row~1 to~2,
row~2 to~3, and row~3 to~4. This is the forbidden transitive
omission of $h_1$. Neither rank-$a$ case is possible.

Finally let $a=2$, so $2\leq b\leq5$, and suppose $d=b$.
All three inequalities are equalities: $M$ has $b-1$ members,
its block image is $S_b$, and its two actual orbits are
natural copies of that block action. Thus $\gen M$ is
faithfully diagonal $S_b$ with its inherited string IP.
It cannot split into nontrivial commuting marked factors:
$S_b$ is directly indecomposable for $2\leq b\leq5$
(its elementary normal-subgroup lists give this immediately).
Hence the members of $M$ occupy one interval. The remaining
generator $x$ is an end; reverse so it is first.

An anchor cannot be at $x$, since $\gen M$ has no fixed
point. It must be at the far endpoint of $M$. Its block $p$
is fixed by that endpoint's omission in $M$, and $x$ fixes
the anchor point. The block image of $x$ centralizes the
near-end omission $Q$ in $M$ and fixes $p$. Lemma~\ref{lem:small}
makes it trivial. Thus $x$ is a binary vector of flips of the
two-point blocks. String commutation makes this vector
$Q$-constant, the anchor makes its coordinate at $p$ zero,
and evenness makes its weight even. The same lemma makes
the whole vector zero, contradicting the nonidentity generator.
This excludes $d=b$ and finishes every small block size.
\end{proof}

\paragraph{Exact scope.}
Together with the accepted degree-at-least-twelve theorem,
the anchored bound now holds in every imprimitive degree.
Only the exact endpoint anchor and intransitive omissions make
the small cases work: parity and string commutation alone do not
eliminate the two degree-eight candidates, but their last
omissions are transitive.
No finite trial is used as a proof premise. The already accepted
$S_5$ maximal-string certificate and full forest/common proof
remain paid dependencies. No new whole-source or formal claim
is made.
\end{document}
